\structural{АНОТАЦІЯ}

Туболев М. Курсова робота на тему «Імітаційна модель системи посадки літаків в аеропорту на основі формального опису стохастичною мережею Петрі».

Текстова частина курсової роботи складається із вступу, п'яти розділів, висновків, списку використаних джерел з \kpistatof{self}{sources} найменувань та \kpistatof{self}{appendices} додатків, які викладено на \kpistatof{self}{pages} сторінках. В тексті роботи оформлено \kpistatof{self}{figures} рисунків, \kpistatof{self}{tables} таблиць.

У роботі поставлено завдання визначити кількість посадкових смуг аеропорту, за якої час окупності додаткових смуг мінімальний. Літаки надходять пуассонівським потоком, займають смугу на 35~хвилин, за тривалого очікування потребують дозаправки, а після 140~хвилин відлітають в інший аеропорт.

Імітаційна модель розроблена на основі формального опису стохастичною мережею Петрі з часовими затримками, у якій кожен літак представлено окремим маркером-об'єктом. Це дозволило точно відтворити індивідуальні таймери дозаправки й відльоту кожного літака. Алгоритм імітації мережі Петрі реалізовано мовою TypeScript у середовищі виконання Bun без бібліотек імітаційного моделювання. Коректність алгоритму підтверджено модульними тестами та верифікацією моделі за вісьмома наборами вхідних даних.

Проведено дослідження перехідного періоду, визначено тривалість прогону (рік модельного часу після доби розгону) і кількість повторних прогонів (51) для довірчої ймовірності 0,95. Повним перебором досліджено від 0 до 10 додаткових смуг.

Результати моделювання свідчать, що за двох смуг аеропорт збитковий: майже кожен літак дозаправляється, дві третини літаків відлітають. Найменший час окупності, близько 0,93~місяця, досягається за семи смуг, тобто за п'яти додаткових; за них жоден літак не відлітає, а середнє очікування скорочується до 9~хвилин.

Зроблено висновок, що аеропорту доцільно збудувати п'ять додаткових смуг. Пробний запуск за інтенсивності 8 і 12~літаків на годину показав, що оптимум зсувається до 5 і 8~смуг відповідно, тож за очікуваного зростання потоку варто передбачити шосту додаткову смугу.

Ключові слова: імітаційна модель, алгоритм імітації, стохастична мережа Петрі, аеропорт, час окупності, TypeScript.
